\documentclass{article}
\usepackage[UTF8]{ctex}
\usepackage{amsmath,amssymb,amsfonts}
\usepackage{enumitem}
\usepackage{geometry}
\geometry{a4paper, margin=1in}
% === 错题本排版优化（不参与题目标识） ===
\setlength{\parindent}{0pt}
\setlength{\parskip}{0.35em plus 0.10em minus 0.08em}
\setlength{\emergencystretch}{3em}
\setlist{itemsep=0.25em, topsep=0.40em, parsep=0pt}
\allowdisplaybreaks[2]
% === 错题本排版优化结束 ===
\begin{document}
\section*{重庆大学2024年高等代数真题}
\section*{55．重庆大学2024年高等代数试卷}
一．计算题．每题8分，共40分．\\

\subsection{题目1}
1．\(a, b\) 满足什么条件时，实系数多项式 \(f(x)=x^{3}+3 a x+b\) 有重因式。

\subsection{题目2}
2．计算 \(n\) 阶行列式

\[
D_{n}=\left|\begin{array}{ccccc}
1+a_{1} & 1 & 1 & \cdots & 1 \\
1 & 1+a_{2} & 1 & \cdots & 1 \\
1 & 1 & 1+a_{3} & \cdots & 1 \\
\vdots & \vdots & \vdots & & \vdots \\
1 & 1 & 1 & \cdots & 1+a_{n}
\end{array}\right| .
\]

其中 \(a_{i} \neq 0(i=1,2, \cdots, n)\) ．

\subsection{题目3}
3．设线性方程组

\[
\left\{\begin{array}{l}
k x_{1}+x_{2}+x_{3}=k-3 \\
x_{1}+k x_{2}+x_{3}=-2 \\
x_{1}+x_{2}+k x_{3}=-2
\end{array}\right.
\]

\begin{itemize}
  \item[（1）] \(k\) 为何值时，方程组无解？
  \item[（2）] \(k\) 为何值时，方程组存在唯一解？
  \item[（3）] \(k\) 为何值时，方程组存在无穷多解？并求出通解。
\end{itemize}

\subsection{题目4}
4．在线性空间 \(P^{3}\) 中，已知 \(\alpha_{1}=(1,1,0), \alpha_{2}=(0,1,1), \alpha_{3}=(0,0,1)\) ，线性变换定义如下：

\[
\mathcal{A} \alpha_{1}=(2,1,1), \mathcal{A} \alpha_{2}=(1,0,1), \mathcal{A} \alpha_{3}=(1,1,2) .
\]

求 \(\mathcal{A}\) 在 \(\beta_{1}=(1,0,0), \beta_{2}=(-1,1,0), \beta_{3}=(0,-1,1)\) 下的矩阵。

二．（10 分）解答如下问题：\\
（1）设不可约多项式 \(p(x)\) 是多项式 \(f(x)\) 的 \(k(k \geq 2)\) 重因式．证明：\(p(x)\) 是 \(f^{\prime}(x)\) 的 \(k-1\) 重因式．并用反例说明反之不成立．\\

\subsection{题目5}
原题第（2）问：求一个三次多项式 \(f(x)\) ，使得 \(f(x)+1\) 能被（x－1）\({ }^{2}\) 整除，\(f(x)-1\) 能被（x＋1）整除．

三。（15分）设矩阵 \(A=\left(\alpha_{1}, \alpha_{2}, \alpha_{3}\right), B=\left(\alpha_{1}, \alpha_{2}, \alpha_{3}, \beta+\alpha_{3}\right)\) ，其中 \(\alpha_{1}, \alpha_{2}, \alpha_{3}, \beta\) 为4维列向量。方程组 \(A X=\beta\) 的通解为 \(X=\xi_{0}+k \xi\) ，其中 \(k\) 为任意常数，\(\xi=\left(\begin{array}{l}1 \\ 1 \\ 1\end{array}\right), \xi_{0}=\left(\begin{array}{l}1 \\ 2 \\ 3\end{array}\right)\) ．\\
（1）求矩阵 \(A, B\) 的秩．\\
（2）求方程组 \(B Y=\alpha_{1}-\alpha_{2}\) 的通解。

四．（15 分）设 \(A, B\) 为 \(n\) 阶矩阵，\(C=\left(\begin{array}{cc}A & O \\ O & B\end{array}\right)\)

\begin{itemize}
  \item[（1）] 当 \(A, B\) 可逆时，求 \(C\) 的逆矩阵和伴随矩阵。
  \item[（2）] 当 \(A, B\) 中至少一个不可逆时，证明 \(C\) 也不可逆．
  \item[（3）] 当 \(A, B\) 中至少一个不可逆时，求 \(C\) 的伴随矩阵．
\end{itemize}

五．（15分）设 \(A\) 是 \(n\) 阶实对称矩阵，证明：\(A\) 可逆当且仅当存在实矩阵 \(B\) ，使得 \(A B+B^{\prime} A\) 正定，其中 \(B^{\prime}\) 是 \(B\) 的转置矩阵。

六．（10 分）设 \(V\) 为 \(n\) 维线性空间， \(1 \leq r<n\) ，证明：\(V\) 有无穷多个 \(r\) 维的子空间．

七．（15 分）设向量 \(\alpha=\left(a_{1}, a_{2}, \cdots, a_{n}\right), \beta=\left(b_{1}, b_{2}, \cdots, b_{n}\right)\) ，令 \(c=\beta \alpha^{\prime}=a_{1} b_{1}+a_{2} b_{2}+\cdots+a_{n} b_{n}, A=\) \(\alpha^{\prime} \beta\) ，其中 \(\alpha^{\prime}\) 是 \(\alpha\) 的转置。

\begin{itemize}
  \item[（1）] 证明：当 \(c \neq 0\) 时，\(A\) 可对角化．
  \item[（2）] 当 \(c=0\) 时，\(A\) 是否可对角化？并说明理由．
\end{itemize}

八．（15 分）设数域 \(P\) 上的 \(n\) 阶方阵 \(A\) 是指数为 \(k\) 的幂零矩阵，即 \(A^{k-1} \neq O, A^{k}=O\) 。其若尔当标准形为

\[
J=\left(\begin{array}{ccc}
J_{1} & & \\
& J_{2} & \\
& & \\
& & J_{s}
\end{array}\right) \text {, 其中 } J_{i}=\left(\begin{array}{cccc}
0 & 1 & & \\
& 0 & \ddots & \\
& & & 1 \\
& & & 0
\end{array}\right)_{n_{i} \times n_{i}}
\]

\begin{itemize}
  \item[（1）] 证明：\(s=n-r(A)\) ．
  \item[（2）] 证明：\(k=\max \left\{n_{1}, n_{2}, \cdots, n_{s}\right\}\) 。
  \item[（3）] 利用（1）和（2）的结论求矩阵 \(A=\left(\begin{array}{ccc}3 & 6 & -15 \\ 1 & 2 & -5 \\ 1 & 2 & -5\end{array}\right)\) 的若尔当标准形．
\end{itemize}

九．（15 分）设二次型 \(f\left(x_{1}, x_{2}, x_{3}\right)\) 的矩阵为 \(A\) ，其在正交变换 \(X=Q Y\) 下的标准形为 \(y_{1}^{2}+y_{2}^{2}\) ，其中\(X=\left(\begin{array}{l}x_{1} \\ x_{2} \\ x_{3}\end{array}\right), Y=\left(\begin{array}{l}y_{1} \\ y_{2} \\ y_{3}\end{array}\right)\) ，且 \(Q\) 的第三列为 \(\left(\begin{array}{c}\frac{\sqrt{2}}{2} \\ 0 \\ \frac{\sqrt{2}}{2}\end{array}\right)\) ．

\begin{itemize}
  \item[（1）] 求 \(A\) 。
  \item[（2）] 证明 \(A+E\) 为正定矩阵，其中 \(E\) 为3阶单位矩阵。
\end{itemize}
\end{document}
